\chapter{Platform Integration}\label{ch:platforms}

This chapter gives the steps to add QuickEIP to a project on each shipped
platform. Table~\ref{tab:platforms} gives the test status of each port.

% -------------------------------------------------------------------------
\section{NetBurner NNDK}\label{sec:netburner}

The port was tested with NNDK 3.5.8 on ARM platforms. The NNDK 3.x socket
API is the same on all platforms.

\begin{procedure}{Add QuickEIP to an NNDK project}
\begin{enumerate}
  \item Put QuickEIP in a subdirectory of the project, for example
        \file{MyProject/QuickEIP}. A git submodule is also correct.
  \item Write the makefile of the project (listing below).
  \item Set the NNDK environment variables.
  \item Run \code{make}.
  \item Find the image at \file{obj/release/MyProject.bin}.
\end{enumerate}
\end{procedure}

\begin{note}
\notelabel The NNDK rules put each object file at
\file{obj/release/}\textit{source path}. A source path that starts with
\file{../} thus puts the object files outside \file{obj/}, next to your
sources. Keep QuickEIP inside the project directory to prevent this.
\end{note}

This is a complete makefile, with QuickEIP in a subdirectory:

\begin{lstlisting}[language=make]
NAME    = MyProject

CPP_SRC += src/main.cpp

QUICKEIP = QuickEIP

C_SRC   += $(QUICKEIP)/src/eip_codec.c \
           $(QUICKEIP)/src/eip_adapter.c \
           $(QUICKEIP)/src/eip_session.c \
           $(QUICKEIP)/src/eip_io.c \
           $(QUICKEIP)/src/eip_scan.c \
           $(QUICKEIP)/src/eip_server.c \
           $(QUICKEIP)/src/eip_adapter_default.c
CPP_SRC += $(QUICKEIP)/ports/netburner/eip_hal_netburner.cpp
NBINCLUDE += -I$(QUICKEIP)/include -I$(QUICKEIP)/src

include $(NNDK_ROOT)/make/boilerplate.mk
\end{lstlisting}

Set the environment and build, for example:

\begin{lstlisting}[language=shell]
export NNDK_ROOT=C:/nburn PLATFORM=MODM7AE70 DEFPLATFORM=MODM7AE70
export PATH=/c/nburn/gcc/bin:/c/nburn/pcbin:/c/nburn/gcc/msys/1.0/bin:$PATH
make
\end{lstlisting}

The first build takes several minutes, because the NNDK compiles its system
library into the \file{obj/} directory of the project. This is normal. For
NBEclipse, add the same files to the project. Add the two include directories
to the C and the C++ compiler settings.

\subsection{Start QuickEIP from UserMain}

Scanner calls block, so run them in a task of their own. Each task needs its
own priority, and a lower number runs first. The adapter task uses
\code{MAIN_PRIO - 1}. Thus \code{MAIN_PRIO + 1} is a safe priority for the
scanner task. \code{UserMain} must then sleep, so that the scanner task can
run.

\begin{lstlisting}[language=C++]
#include <init.h>
#include <nbrtos.h>
#include <string.h>

#include "eip.h"

const char *AppName = "MyProject";

static void ScannerTask(void *)
{
    for (;;) {
        /* eip_session_open, eip_tag_read, and so on (chapter 6) */
        OSTimeDly(TICKS_PER_SECOND);
    }
}

void UserMain(void *)
{
    EipAdapterConfig cfg;

    init();
    WaitForActiveNetwork(TICKS_PER_SECOND * 10);

    memset(&cfg, 0, sizeof cfg);
    /* fill in cfg as in chapter 4 */
    eip_adapter_start(&cfg);
    OSSimpleTaskCreatewName(ScannerTask, MAIN_PRIO + 1, "EipScanner");

    for (;;) {
        /* eip_adapter_read and eip_adapter_write */
        OSTimeDly(1);
    }
}
\end{lstlisting}

\subsection{Platform Notes}

\begin{description}[font=\sffamily\bfseries\color{steel}]
  \item[Clock resolution.] \code{eip_hal_millis} follows \code{TimeTick}: 50~ms
        at the default 20~Hz. The adapter never sets a Class~1 timeout shorter
        than two ticks (100~ms), and \TO{} packets go out on tick boundaries.
        For RPIs shorter than approximately 50~ms, increase
        \code{TICKS_PER_SECOND} in the platform configuration.
  \item[Task.] The adapter task runs at \code{MAIN_PRIO - 1} by default. Set
        \code{EIP_ADAPTER_TASK_PRIORITY} to change it. The task uses the
        default NNDK task stack.
  \item[Lock.] \code{eip_hal_lock} is \code{OSLock()}, which stops task
        switches for a short time. It is never held during a socket call.
  \item[UDP sockets.] The port uses \code{CreateRxTxUdpSocket4} with an
        explicit local port. When the library asks for ``any port,'' the port
        selects one in 50000--50063.
  \item[TCP options.] The port does not set TCP\_NODELAY or keep-alive. A slot
        of a lost peer comes back through eviction after
        \code{EIP_ADAPTER_EVICT_MS}.
  \item[Web UI.] \file{examples/netburner_webui} is a test web page. It is
        application code, not part of the stack. It has no \code{main} and no
        makefile, so add it to an NNDK application.
\end{description}

% -------------------------------------------------------------------------
\section{ESP-IDF (ESP32)}\label{sec:espidf}

The port was tested with ESP-IDF 5.5.1 on the esp32 target. All ESP-IDF 5.x
versions should operate.

\subsection{Add the Component}

The root directory of QuickEIP is an IDF component. The component takes the
name of its directory. Use the same spelling in all locations. The simpler
method is a top-level \file{CMakeLists.txt} that names the directory. The
path is relative to the project. Here QuickEIP is beside it:

\begin{lstlisting}[language=cmake]
cmake_minimum_required(VERSION 3.16)
list(APPEND EXTRA_COMPONENT_DIRS "${CMAKE_CURRENT_LIST_DIR}/../QuickEIP")
include($ENV{IDF_PATH}/tools/cmake/project.cmake)
project(myapp)
\end{lstlisting}

With the component manager, put the dependency in \file{main/idf_component.yml}.
The path is relative to \file{main/}:

\begin{lstlisting}[language=yaml]
dependencies:
  QuickEIP:
    path: ../../QuickEIP
\end{lstlisting}

\begin{caution}
The dependency key must be the same as the directory name. If it is not, the
build fails with \code{Failed to resolve component ... unknown name}. A change
to \code{REQUIRES} does not correct this error.
\end{caution}

With the two methods, \file{main/CMakeLists.txt} names the component and the
network components that you use:

\begin{lstlisting}[language=cmake]
idf_component_register(SRCS "main.c"
                       INCLUDE_DIRS "."
                       REQUIRES QuickEIP esp_eth esp_netif esp_event)
\end{lstlisting}

Build with \code{idf.py set-target esp32} and then \code{idf.py build}.

\subsection{Start Ethernet}

Use the driver code in \file{examples/ethernet/basic} of ESP-IDF as a model.

\begin{warning}
You must know the PHY, the pins, and the RMII clock mode of your board. A
wrong value compiles without an error, but the link never comes up. Get these
values from the schematic of the board. Do not guess them.
\end{warning}

This example is for the internal EMAC of the ESP32 with a LAN87xx PHY:

\begin{lstlisting}
#include "esp_eth.h"
#include "esp_event.h"
#include "esp_netif.h"

static esp_eth_handle_t s_eth;
static esp_netif_t *s_netif;
static volatile int s_link_up;

static void eth_event(void *arg, esp_event_base_t base, int32_t id, void *data)
{
    (void)arg; (void)base; (void)data;
    if (id == ETHERNET_EVENT_CONNECTED) s_link_up = 1;
    else if (id == ETHERNET_EVENT_DISCONNECTED) s_link_up = 0;
}

static void eth_start(void)
{
    eth_mac_config_t mac_cfg = ETH_MAC_DEFAULT_CONFIG();
    eth_esp32_emac_config_t emac_cfg = ETH_ESP32_EMAC_DEFAULT_CONFIG(); /* set smi_gpio and clock_config */
    eth_phy_config_t phy_cfg = ETH_PHY_DEFAULT_CONFIG();                /* and your PHY address */
    esp_eth_mac_t *mac = esp_eth_mac_new_esp32(&emac_cfg, &mac_cfg);
    esp_eth_phy_t *phy = esp_eth_phy_new_lan87xx(&phy_cfg);
    esp_eth_config_t eth_cfg = ETH_DEFAULT_CONFIG(mac, phy);
    esp_netif_config_t netif_cfg = ESP_NETIF_DEFAULT_ETH();

    ESP_ERROR_CHECK(esp_eth_driver_install(&eth_cfg, &s_eth));
    ESP_ERROR_CHECK(esp_netif_init());
    ESP_ERROR_CHECK(esp_event_loop_create_default());
    s_netif = esp_netif_new(&netif_cfg);
    ESP_ERROR_CHECK(esp_netif_attach(s_netif, esp_eth_new_netif_glue(s_eth)));
    ESP_ERROR_CHECK(esp_event_handler_register(ETH_EVENT, ESP_EVENT_ANY_ID, eth_event, NULL));
    ESP_ERROR_CHECK(esp_eth_start(s_eth));
}
\end{lstlisting}

\subsection{Report the Link}

The esp\_netif layer does not report speed or duplex. Read them from the
driver in a \code{link_info} callback. Return 1 while the link is up, and 0 to
report it down:

\begin{lstlisting}
static int link_info(void *ctx, uint32_t *mbps, int *full, int *autoneg)
{
    eth_speed_t speed;
    eth_duplex_t duplex;
    (void)ctx;
    if (!s_link_up) return 0;
    if (esp_eth_ioctl(s_eth, ETH_CMD_G_SPEED, &speed) != ESP_OK) return 0;
    if (esp_eth_ioctl(s_eth, ETH_CMD_G_DUPLEX_MODE, &duplex) != ESP_OK) return 0;
    *mbps = (speed == ETH_SPEED_100M) ? 100 : 10;
    *full = (duplex == ETH_DUPLEX_FULL);
    *autoneg = 1;
    return 1;
}
\end{lstlisting}

\subsection{Start the Adapter}

\begin{procedure}{Start the adapter on ESP-IDF}
\begin{enumerate}
  \item Call \code{eth_start}.
  \item Set all fields of \code{cfg} (chapter~\ref{ch:adapter}).
  \item Set \code{cfg.link_info} to your callback.
  \item Optional: call \code{eip_esp_set_netif} to select the interfaces.
  \item Call \code{eip_adapter_start}.
\end{enumerate}
\end{procedure}

\begin{lstlisting}
#include "eip.h"
#include "eip_esp_idf.h"

eth_start();
cfg.link_info = link_info;       /* with the rest of cfg, from chapter 4 */
eip_esp_set_netif(&s_netif, 1);  /* optional: default is ETH_DEF, WIFI_STA_DEF, WIFI_AP_DEF */
eip_adapter_start(&cfg);
\end{lstlisting}

Call \code{eip_esp_set_netif} before \code{eip_adapter_start}. The function
copies the list, and the netifs must stay valid while the adapter runs.

\subsection{Options}

Put the options in compile definitions for all components. In the top-level
\file{CMakeLists.txt}, set them after the line that includes
\file{project.cmake} and before \code{project()}:

\begin{lstlisting}[language=cmake]
idf_build_set_property(COMPILE_DEFINITIONS "EIP_ESP_TASK_PRIORITY=12" APPEND)
idf_build_set_property(COMPILE_DEFINITIONS "EIP_ESP_TASK_WDT=1" APPEND)
\end{lstlisting}

\begin{table}[htbp]
\centering
\caption{ESP-IDF port options}\label{tab:espopts}
\begin{tabularx}{\linewidth}{@{}P{0.30\linewidth}P{0.17\linewidth}L@{}}
\toprule
\textbf{Define} & \textbf{Default} & \textbf{Effect} \\
\midrule
\code{EIP_ESP_TASK_STACK} & 6144 & Adapter task stack in bytes (or \code{EIP_ADAPTER_TASK_STACK}). \\
\code{EIP_ESP_TASK_PRIORITY} & 12 & Adapter task priority (or \code{EIP_ADAPTER_TASK_PRIORITY}). \\
\code{EIP_ESP_TASK_CORE} & \code{tskNO_AFFINITY} & Pins the adapter task to a core. \\
\code{EIP_ESP_TASK_WDT} & 0 & 1 adds the adapter task to the task watchdog. \\
\bottomrule
\end{tabularx}
\end{table}

\subsection{Platform Notes}

\begin{itemize}
  \item Use wired Ethernet for Class~1 in production. Wi-Fi latency peaks are
        frequently longer than a timeout of 8 $\times$ RPI.
  \item The adapter task waits in \code{select()}, so it starts as soon as a
        packet arrives. While idle, it still does a pass at least each
        \code{EIP_ADAPTER_IDLE_WAIT_MS} (5~ms).
  \item Arduino-ESP32 runs on ESP-IDF and can use the same files. Make a
        library directory whose \file{src/} contains \file{include/*},
        \file{src/*}, \file{ports/bsd_sockets/*}, and \file{ports/esp-idf/*}.
        Do not include the other ports, because Arduino compiles all files
        that it finds. This setup has not been tested.
  \item The example \file{examples/esp_idf_adapter} uses
        \code{protocol_examples_common}. For that example only, run
        \code{idf.py menuconfig} and set the network under
        \menu{Example Connection Configuration}.
\end{itemize}

% -------------------------------------------------------------------------
\section{Teknic ClearCore}\label{sec:clearcore}

\file{examples/clearcore_adapter} compiles and links with the arm-gcc 6.3 of
Microchip Studio, against the Enhanced ClearCore Library, a fork of ClearCore
library 1.7. The adapter has run on a ClearCore with a Micro850 PLC. The
example itself has not been run.

ClearCore runs on bare metal (no RTOS), on the raw API of lwIP with
\code{NO_SYS=1}. QuickEIP uses lwIP directly through \file{ports/lwip_raw}. It
does not use the \code{EthernetTcp*} and \code{EthernetUdp} classes. Those
classes have a 600-byte TCP buffer and no broadcast send, and EtherNet/IP
needs more.

\begin{procedure}{Set up a ClearCore project}
\begin{enumerate}
  \item Copy \file{ProjectTemplate}, \file{libClearCore}, and \file{LwIP}
        from the ClearCore library directory to a location where you can
        write.
  \item Put your QuickEIP directory beside them. The four directories must be
        at the same level, because the paths below are relative.
  \item Open \file{ProjectTemplate.atsln}.
  \item Make sure that the three projects load.
  \item Add the files with \menu{Add $\to$ Existing Item $\to$ Add As Link}:
        \file{src/*.c} (all files), \file{ports/lwip_raw/eip_hal_lwip_raw.c},
        and \file{ports/clearcore/eip_hal_clearcore.cpp}.
  \item Open \menu{Project Properties $\to$ Toolchain}. Do steps 7--10 in the
        Debug and the Release configurations.
  \item In \menu{ARM/GNU C Compiler $\to$ Directories} and \menu{ARM/GNU C++
        Compiler $\to$ Directories}, add \file{../../QuickEIP/include} and
        \file{../../QuickEIP/src}.
  \item In \menu{ARM/GNU C Compiler $\to$ Directories} only, add
        \file{../../LwIP/LwIP/src/include} and
        \file{../../LwIP/LwIP/port/include}.
  \item In \menu{ARM/GNU C Compiler $\to$ Symbols} and \menu{ARM/GNU C++
        Compiler $\to$ Symbols}, add \code{EIP_ADAPTER_MAX_CLIENTS=2}.
  \item Keep the symbols \code{DEBUG} and \code{NDEBUG} and the flag
        \code{-std=gnu11} that are in the template.
  \item Replace \file{main.cpp} of the template with your main loop.
\end{enumerate}
\end{procedure}

\begin{note}
\notelabel The template sets the lwIP include paths for C++ only. The lwIP
port is a C file. Without step 8, the build stops at \file{lwip/ip_addr.h}.
\end{note}

\begin{caution}
Add \code{EIP_ADAPTER_MAX_CLIENTS=2} to the C and the C++ symbols. It changes
the size of structures, so the two compilers must agree. The default of 4 is
too many for this lwIP build.
\end{caution}

Paths are relative to the \file{Debug} or \file{Release} output directory.
To check a compile without Studio, run \code{arm-none-eabi-gcc} or
\code{arm-none-eabi-g++} with \code{-c} from \file{ProjectTemplate/Debug} and
these flags. The toolchain is in
\file{Atmel/Studio/7.0/toolchain/arm/arm-gnu-toolchain/bin}, and
\textit{Packs} is \file{Atmel/Studio/7.0/Packs}:

\begin{lstlisting}[language=shell]
-mthumb -mcpu=cortex-m4 -mfloat-abi=hard -mfpu=fpv4-sp-d16
-D__SAME53N19A__ -DEIP_ADAPTER_MAX_CLIENTS=2
-I../../QuickEIP/include -I../../QuickEIP/src -I../../QuickEIP/ports/clearcore
-I../../libClearCore/inc -I../../LwIP/LwIP/src/include -I../../LwIP/LwIP/port/include
-I<Packs>/atmel/SAME53_DFP/1.1.118/include
-I<Packs>/arm/CMSIS/4.5.0/CMSIS/Include
-I<Packs>/arm/CMSIS/5.4.0/CMSIS/Core/Include
\end{lstlisting}

\subsection{The Main Loop}

The full example is in \file{examples/clearcore_adapter/main.cpp}:

\begin{lstlisting}[language=C++]
#include "ClearCore.h"
#include "eip.h"

while (!EthernetMgr.PhyLinkActive()) {                 // wait for the cable
    Delay_ms(100);
}
EthernetMgr.Setup();
if (!EthernetMgr.DhcpBegin()) {
    EthernetMgr.LocalIp(IpAddress(192, 168, 1, 50));   // fallback static address
}
if (eip_adapter_start(&cfg) != EIP_OK) {               // poll mode: starts no task
    while (true) { EthernetMgr.Refresh(); }            // stop: outputs stay off
}

while (true) {
    eip_adapter_poll();             // every pass
    EthernetMgr.Refresh();
    // motion, I/O, eip_adapter_read / eip_adapter_write ...
}
\end{lstlisting}

\subsection{Platform Notes}

\begin{description}[font=\sffamily\bfseries\color{steel}]
  \item[Call \code{eip_adapter_poll} frequently.] \TO{} packets go out from
        inside the poll call. Thus the loop time becomes RPI jitter. Keep the
        loop time much shorter than your RPI.
  \item[The lwIP build is small.] It has a 4~KB heap (\code{MEM_SIZE}), five
        TCP PCBs, and eight listening PCBs. The adapter uses one listening PCB
        plus one PCB for each explicit session. A scanner session uses one
        more. This is why you set \code{EIP_ADAPTER_MAX_CLIENTS} to 2 or 3. For
        more, compile the LwIP project again with a larger \code{MEM_SIZE} and
        \code{MEMP_NUM_TCP_PCB}. Its \file{lwipopts.h} uses \code{\#ifndef},
        so \code{-D} operates.
  \item[Blocking scanner calls.] They call \code{EthernetMgr.Refresh()} while
        they wait, and motion continues from interrupts. But your main loop
        does not run until the call returns. Keep \code{timeout_ms} short.
  \item[Interrupts.] Do not call QuickEIP from an interrupt handler. The lwIP
        callbacks run inside \code{EthernetMgr.Refresh()} in the main context.
        This is why the port needs no lock.
  \item[Socket table.] \code{EIP_LWIP_MAX_SOCKETS} (12),
        \code{EIP_LWIP_UDP_QUEUE} (8), and \code{EIP_LWIP_ACCEPT_QUEUE} (4) set
        the size of the raw-API socket table.
  \item[Link report.] The port reports 100~Mbit/s full duplex while the PHY
        link is up.
  \item[Multicast \TO.] The port sets TTL 1 only in lwIP builds with
        \code{LWIP_MULTICAST_TX_OPTIONS}. The standard build does not have
        this option. Thus use a point-to-point connection.
  \item[Host name.] lwIP on ClearCore has no host name. Call
        \code{eip_clearcore_set_host_name} at start-up, and again from
        \code{apply_interface} when a scanner sets attribute 6.
\end{description}

\begin{warning}
On ClearCore, the main loop also controls motion. A blocking scanner call
stops the main loop until it returns. Do not make scanner calls while motion
needs service from the main loop. Keep \code{timeout_ms} short.
\end{warning}

The same three groups of files should also operate as an Arduino library with
the ClearCore wrapper. This has not been tested.

% -------------------------------------------------------------------------
\section{Linux, macOS, and WSL}\label{sec:posix}

Section~\ref{sec:hostbuild} gives the build procedure. The build makes
\file{libquickeip.a} (the core and the POSIX port), the tests, and
\code{eiptool}. Appendix~\ref{app:tools} describes \code{eiptool}.

\begin{itemize}
  \item Set \code{EIP_NETIF=eth1} to select the interface that the adapter
        reports and that the scan uses first. This applies to all programs that
        use the POSIX port. The port ignores an unknown name and uses the
        default interface.
  \item Ports 44818 and 2222 are above 1024, so you do not need root access.
        Another EtherNet/IP program on the same host blocks them.
  \item Under WSL2, the host is behind NAT. A scan to the LAN operates, but
        PLCs on your LAN cannot get to a WSL adapter. Use mirrored networking.
\end{itemize}

% -------------------------------------------------------------------------
\section{Windows}\label{sec:windows}

The port was tested with Visual Studio 2022 on Windows 11, as x64 and Win32.
The port is \file{ports/windows/eip_hal_windows.c}. It needs Windows 7 or
later. Section~\ref{sec:hostbuild} gives the build procedure.

\begin{procedure}{Add QuickEIP to a Windows project}
\begin{enumerate}
  \item Compile \file{src/*.c} and \file{ports/windows/eip_hal_windows.c}.
  \item Add \file{include/} and \file{src/} to the include path.
  \item Link \code{ws2_32}, \code{iphlpapi}, and \code{winmm}.
\end{enumerate}
\end{procedure}

\begin{table}[htbp]
\centering
\caption{Windows port options}\label{tab:winopts}
\small
\begin{tabularx}{\linewidth}{@{}P{0.32\linewidth}P{0.10\linewidth}L@{}}
\toprule
\textbf{Define} & \textbf{Default} & \textbf{Effect} \\
\midrule
\code{EIP_WIN_POLL_MODE} & Off & Starts no adapter thread. Your loop calls \code{eip_adapter_poll}. \\
\code{EIP_ADAPTER_TASK_PRIORITY} & 0 & A \code{THREAD_PRIORITY_*} value for the adapter thread. 0 selects \code{THREAD_PRIORITY_HIGHEST}. \\
\code{EIP_ADAPTER_TASK_STACK} & 0 & Adapter thread stack in bytes. 0 is the system default. \\
\code{EIP_WIN_MAX_SOCKETS} & 64 & Sockets open at the same time. \\
\code{EIP_WIN_MAX_INTERFACES} & 16 & Interfaces reported. \\
\code{EIP_WIN_NETIF_REFRESH_MS} & 2000 & Interval at which the port reads the interface list again, outside the adapter thread. \\
\bottomrule
\end{tabularx}
\end{table}

\subsection{Platform Notes}

\begin{description}[font=\sffamily\bfseries\color{steel}]
  \item[Select the interface on a PC with more than one.] The adapter reports
        interface 0 as its address, in ListIdentity and in the \OT{} socket
        address of a Forward\_Open. Interface 0 is the first one that is up
        and has a default gateway. Frequently, this is not the machine
        network. Set \code{EIP_NETIF} to the friendly name or the IPv4 address
        of the adapter interface, for example \code{EIP_NETIF="Ethernet 2"} or
        \code{EIP_NETIF=10.168.41.1}. The port keeps a named interface also
        while its cable is disconnected, and reports the link as down. Scanner
        calls do not need this setting.
  \item[Firewall.] An adapter needs inbound TCP 44818, UDP 44818, and UDP
        2222 for your program. A scanner that runs Class~1 needs inbound UDP
        on its local port. Use \code{Get-NetConnectionProfile} to find the
        profile of the network.
  \item[A new executable can lose its first replies.] Third-party security
        software can stop inbound UDP for the first seconds that a new
        program runs, also with Windows Firewall off. On a test PC with AVG,
        each new binary failed its first \code{eiptool identity}, and then
        answered in approximately 45~ms. If a new build times out, or a scan
        finds too few devices, run it again before you look for a network
        fault.
  \item[Busy ports.] Other EtherNet/IP software on the PC, such as RSLinx,
        can hold 44818 or 2222. Each bind is exclusive. The adapter then
        cannot open the ports, and \code{listening} stays 0.
  \item[Timing.] The port sets a 1~ms system timer period for the process
        from its first call. This uses more power on a laptop. On the test PC,
        a 10~ms RPI had a mean of 10.00~ms and a maximum of 11.83~ms over
        \num{1500} packets, with the two ends on the same host.
  \item[Link report.] Windows gives the speed. It does not give duplex or
        autonegotiation, so the port reports full duplex and autonegotiated
        while the link is up.
  \item[Not tested.] Multicast \TO, and a real PLC connected to a Windows
        adapter.
\end{description}

\begin{warning}
Windows is not a real-time operating system. A stall of approximately 100~ms
was measured on a bench PC. Set \code{timeout_floor_ms} to 500 or more. Do
not use a Windows adapter for a safety function.
\end{warning}
